\section{Analysis}
\label{sec:analysis}

The hypotheses ask whether an effect exists; the per-decision logs let us ask where it comes from. The analyses below are fixed before the confirmatory runs, and every quantity in them is computed from \texttt{steps.jsonl} and \texttt{episodes.jsonl}.

\subsection{Format failures and decision errors}
\label{sec:analysis:errors}

Each decision under \texttt{prompt\_generate} falls into one of four classes. A \emph{format failure} produced no legal label and was replaced by a random action; these are the decisions logged as illegal. A \emph{lenient parse} yielded a legal label only through one of the parser's fallback rules, a bare label on a late line or the last label mentioned anywhere, rather than through an \texttt{Action:} line; we recover this class by re-running the parser's stages on the logged raw output. A \emph{decision error} is a legal action that is not optimal, and the remaining decisions are optimal. We judge optimality against the set of optimal actions, not the single action the oracle logged. Gridworld often has several shortest paths and tic-tac-toe several equally good cells, and every episode is deterministic given its seed and the actions taken, so we replay each episode, recover every visited state and compute its full optimal set.

The four rates, by model, method and task, show how much of the gap between small and large models is formatting. \result{[Share of the generation-minus-oracle gap attributable to format failures, by size and task]}. The comparison with scoring can be made decision by decision on loan approval, where the applicants do not depend on earlier actions, so every policy sees the same states; on the other tasks only the first decision of each episode is shared. On those shared states, generation and scoring choose the same legal action in \result{[agreement rate by size]} of the decisions where generation produced one. If the two agree on most such decisions, the benefit of scoring in H4 is mostly the removal of format failures; if they do not, scoring also changes which legal action the model prefers.

Chain-of-thought and thinking outputs, which the reasoning block records on gridworld, tic-tac-toe and support triage, let us look for the knowing-doing gap of \citet{schmied2026greedy}. For each such decision we take the reasoning, either the text before the final \texttt{Action:} line or the \texttt{<think>} block, and check whether it names an optimal action or states the fact that makes an action optimal: an immediate win or a forced block in tic-tac-toe, the shorter way around a wall in gridworld, the rule of the written policy that applies to the ticket in support triage. We then check whether the final action differs. Automatic matching proposes candidates and a stratified sample is annotated by hand \result{[sample size and agreement between automatic and manual labels]}. Only the explicit \texttt{Action:} line counts as the final action here, because the parser's last-mention rule could otherwise take the action from the reasoning itself. \result{[Rate of decisions with a correct rationale and a wrong action, by size and task]}.

On the bandit we separate the failure modes named in prior work. Greediness is the share of arms never pulled \citep{schmied2026greedy}; frequency bias is the probability of repeating the arm chosen most often so far; and a failure to exploit is choosing, late in an episode, an arm other than the one with the highest empirical mean among arms pulled at least $k$ times \citep{harris2026explore}, with $k$ fixed before the confirmatory runs \todo{choose $k$ and add it to the pre-registration}. The suboptimal-pull fraction and the share of the last ten pulls on the best arm are logged per episode. \result{[The three rates by size and method]}.

Errors also have a direction. In blackjack, standing where the table hits is conservative and hitting where it stands is aggressive; we tabulate both by player total, soft or hard hand, and dealer card, to see whether models are systematically risk-averse, as \citet{jia2024decision} found language models to be in multiple-choice-list experiments \result{[conservative and aggressive error rates by size]}. In loan approval the analogous quantity is the approval rate relative to the oracle's \result{[approval rates]}. Under the sign flip we regenerate each applicant from its seed and ask how many errors follow the usual meaning of the credit score, that is, how many decisions match what the unflipped default model would prescribe \result{[share of prior-following errors by size and method]}.

\subsection{Quality versus cost}
\label{sec:analysis:cost}

Latency per decision mixes model size, method and hardware, so we report it next to token counts, which do not depend on hardware. All our latencies come from one laptop, and those of the main sweep from one runtime, PyTorch on Metal, which is slower than MLX or llama.cpp on the same machine (Section~\ref{sec:setup:hw}); ratios between cells travel better than their absolute values. Scoring needs one forward pass over the prompt and the candidate labels, generation adds one decoding step per generated token, the probe needs one forward pass and a logistic regression, and LoRA adds a one-off training cost (\texttt{prepare\_time\_s} in \texttt{summary.json}) to inference by scoring, a cost large enough on the laptop that we stopped LoRA at 4B. Suppose a small adapted model with latency $\ell_{\text{small}}$ per decision matches a larger zero-shot model with latency $\ell_{\text{large}}$. Its training time $T_{\text{train}}$ is repaid after
\begin{equation}
n^\star = \frac{T_{\text{train}}}{\ell_{\text{large}}-\ell_{\text{small}}}
\end{equation}
decisions, \result{[$n^\star$ for the H3 pairs that match]}. The frontier in Figure~\ref{fig:pareto} leaves this training cost out; including it moves \result{[which cells drop off the frontier]}. For H7c we divide \norm{} by latency per decision. A negative \norm{} makes that ratio hard to read, so we compare ratios only between cells whose scores are positive \todo{confirm this reading of the pre-registered H7c criterion before the tag}.

\subsection{What the feature probe controls for}
\label{sec:analysis:feature}

The feature probe is a control, but what it controls for depends on what the feature vector $\phi$ contains, and the seven tasks differ in this respect.

On loan approval the default model in Equation~\eqref{eq:loan} is a logistic function of the displayed fields, linear apart from the logarithmic loan-to-income term and the cap on years of employment, and the oracle thresholds it. A linear head on $\phi$ can therefore come close to the oracle given enough training states, and it sets a high bar: a language-model method that does not clear it adds nothing beyond the numbers it was shown. On the shifted variants the feature probe, like every adapted method, is trained on the unshifted task. Under the sign flip it has learned the wrong sign for the credit score and should fall far below its in-distribution score; under covariate shift the default model is unchanged, but the probe must extrapolate to lower incomes and credit scores. Its drop under each shift is the reference against which H3b's drops for language models are read.

On the contextual bandit $\phi$ is the current context alone, and the weights $w_a$ are redrawn for every episode. For any fixed function of the context, each option is then equally likely to be the best one, so the expected normalized score of the feature probe is exactly zero and the control measures nothing on this task. A language-model method that does better must be using the history in its prompt, which is the in-context learning the task exists to test.

On the bandit $\phi$ holds pull frequencies and empirical means, and the training labels come from a clairvoyant oracle. Under the behavior policy the oracle's arm is also the most-pulled arm, so a cloned policy can learn to repeat whichever arm it has pulled most, which is pure exploitation \citep{laskin2023incontext,nie2025evolve}. The same applies to \texttt{probe} and \texttt{lora\_sft}, whose prompts show the same counts, and to the few-shot demonstrations, whose states come from oracle rollouts, so that the only arm ever pulled in them is the oracle's. Low bandit scores for adapted methods therefore say more about cloning a clairvoyant teacher than about the model.

On gridworld, tic-tac-toe and blackjack the oracle is a non-linear function of $\phi$: a shortest path around walls, a game-tree search, and a lookup table whose stand region is not monotone in the dealer's card. A linear head on raw features is a weak control there. The gap between \texttt{probe} and \texttt{feature\_probe} measures how much of the decision the model's representation makes linearly decodable, which is not the same as the model using it \citep{belinkov2022probing}. The gap between \texttt{probe} and \texttt{prompt\_score} on the same model measures how much the model knows but does not show \citep{orgad2025know}, and the comparison of middle and last layers tests whether that information peaks before the output \citep{skean2025layer}. \result{[Probe minus feature probe and probe minus scoring, by task and size]}. The training statistics in \texttt{summary.json}, probe training accuracy and LoRA loss before and after training, separate a method that failed to fit its training states from one that fitted them and failed to generalize \result{[cells with low training accuracy]}.

\subsection{Oracle agreement and return}
\label{sec:analysis:agreement}

Oracle agreement, the share of decisions in which the policy's own legal action equals the logged oracle action, is easy to read but has two flaws, which is why return is the primary outcome. It is a lower bound on the share of optimal decisions whenever several actions are optimal, because the oracle breaks ties by a fixed rule. In tic-tac-toe against the random opponent the oracle does not prefer quicker wins: in the state shown in Appendix~\ref{app:prompt:tictactoe}, four cells win with certainty, including an immediate win at cell 7, and the oracle picks cell 5. On the two bandits it also measures agreement with a clairvoyant oracle that no learner can match early in an episode; UCB1 agrees with it in \result{[share]} of its decisions. Fallback actions after format failures never count as agreement. For each task we report the rank correlation between agreement and \norm{} across cells \result{[Kendall's $\tau$ by task]} and look at the cells where the two disagree. High agreement with low return points to errors at decisive states, such as a missed block in tic-tac-toe; low agreement with high return points to alternative optimal actions, which the set-valued optimality of Section~\ref{sec:analysis:errors} resolves.
